\section*{Chapter \ref{ch:hf:auctions-and-the-close} --- Auctions and the Close}

\begin{solution}{exo:hf:auctions-and-the-close:1}
Supply reaches 30 shares at 102 (10 at 101 plus 20 at 102): the auction closes at 102.
\end{solution}

\begin{solution}{exo:hf:auctions-and-the-close:2}
$9-0.15\times14=6.9$ cents a share; \$8\,280 on 120\,000 shares.
\end{solution}

\begin{solution}{exo:hf:auctions-and-the-close:3}
Before 3:55 p.m.; every second from 3:55 p.m.
\end{solution}

\begin{solution}{exo:hf:auctions-and-the-close:4}
Each extra share sold lowers the close for all the shares sold; the marginal share's profit falls to zero before the imbalance is filled
(120\,000 of 200\,000 here), and the auction's other sellers fill the rest at a higher price.
\end{solution}

\begin{solution}{exo:hf:auctions-and-the-close:5}
The profit is the close less tomorrow's value; when half the 14-cent move stays, tomorrow's value is 7 cents higher and the provider also offers
less (80\,000 shares), so its close is higher but its margin smaller: 4 cents a share instead of 8.
\end{solution}

\begin{solution}{exo:hf:auctions-and-the-close:6}
If late sell orders shrink the imbalance, an offer at the reference would set a close too low for the provider to profit; a limit 4 cents above
keeps the provider out of those auctions, and costs little when the imbalance holds: 6.0 cents a share against 5.5.
\end{solution}

\begin{solution}{exo:hf:auctions-and-the-close:7}
The imbalance moves the close 10 cents instead of 14; the provider's best offer is 100\,000 shares at the reference for 5.5 cents a share
(against 6.9): a deeper book leaves less to earn.
\end{solution}

\begin{solution}{exo:hf:auctions-and-the-close:8}
The average hides the days when the imbalance carries news: then most of the move stays and a full-size offset loses. The provider must size to
the expected profit, which falls as the share that stays rises, and cut size when the imbalance coincides with news.
\end{solution}

\begin{solution}{pb:hf:auctions-and-the-close:1}
\begin{enumerate}
\item Index funds, benchmarked managers and derivatives settle at the closing price.
\item 7.5\% of daily volume in 2018 (3.1\% in 2010); closing prices typically at the pre-close bid or ask; lower impact than continuous trading;
deviations reverting quickly and almost completely on average.
\item Publications every 10 seconds from 3:50, every second from 3:55; market-on-close orders before 3:55; limit-on-close before 3:58;
modifications before 3:50.
\item The reference price, paired shares, imbalance shares and side, and near and far indicative prices in the last minutes.
\item See \cref{def:hf:auctions-and-the-close:offset}.
\item See \cref{def:hf:auctions-and-the-close:impact}; 5, 10, 14 and 20 cents for 25\,000, 100\,000, 200\,000 and 400\,000 shares.
\item $q\,(p-P_0-\pi\Delta(I))$ for $q$ shares sold at the close $p$.
\item The close falls with every share sold.
\item 120\,000 shares at the reference; the close moves 9 cents instead of 14; 6.9 cents a share, \$8\,280.
\item It rises less than proportionally with the imbalance (4.5, 6.9 and 10 cents at 100\,000, 200\,000 and 400\,000) and falls with $\pi$ (8, 6.9 and 4
cents at 200\,000 for 0, 15\% and 50\%).
\item From 6.9 to 5.5 cents a share; a 4-cent premium recovers 6.0.
\item From 14 cents to 9.
\item With 15\% of the move staying: 2.25, 2.95, 4.5, 6.9, 8.45 and 10 cents a share for imbalances of 25\,000 to 400\,000 shares; 85\% of the
closing move reverts by assumption, and the profit is most sensitive to that assumption.
\item Regress the next open's price against the close's deviation from the pre-close mid, by stock and by type of day.
\item Index futures at settlement for the market part; the idiosyncratic part only by size.
\item Less information, stale reference prices: smaller sizes, larger premiums.
\item It lowers it, by moving the close back toward the reference and taking part of the imbalance.
\item It makes the close nearer to the value the next day confirms.
\item Closing-auction provision on news days.
\item Immediacy at the close, paid for by the part of the move that reverts.
\end{enumerate}
\end{solution}

\begin{solution}{iq:hf:auctions-and-the-close:1}
The imbalance relative to the stock's closing volume and depth, the reference price against the continuous book, any news, index events, the
history of the stock's closing reversals, and the firm's overnight risk budget.
\iqlookfor{size relative to depth, news, reversal history.}
\end{solution}

\begin{solution}{iq:hf:auctions-and-the-close:2}
For each close, the deviation of the closing price from the pre-close mid, and the next morning's (or next day's) price against the close; regress
the later move on the deviation, by stock and period.
\iqlookfor{deviation and subsequent move.}
\end{solution}

\begin{solution}{iq:hf:auctions-and-the-close:3}
Maximise executed volume, then minimise surplus, then follow market pressure, then the reference price, with a fixed rule for exact ties (the
lower price, or the one nearest the reference).
\iqlookfor{the hierarchy of rules.}
\end{solution}

\begin{solution}{iq:hf:auctions-and-the-close:4}
The sum of 800 overnight positions: a market exposure (hedged with futures) and 800 idiosyncratic gaps; limit it by a firm-wide overnight value at
risk and a per-stock size relative to the imbalance and the stock's volatility.
\iqlookfor{aggregation and limits.}
\end{solution}

\begin{solution}{iq:hf:auctions-and-the-close:5}
Index funds must trade the rebalance at the close whatever the price, so imbalances are large and predictable; providers who offset them earn the
reversal.
\iqlookfor{forced demand and its providers.}
\end{solution}

\begin{solution}{iq:hf:auctions-and-the-close:6}
The close is $P_0+\lambda(I-q)$, so profit is $q(\lambda(I-q)-\pi\lambda I)$, maximised at $q^\ast=I(1-\pi)/2$.
\iqlookfor{the first-order condition.}
\end{solution}
